\chapter*{使用说明}
\addcontentsline{toc}{chapter}{使用说明}
\markboth{使用说明}{}

\section*{一、为什么课后习题是 818 复习的第一优先级}

《西华大学 818 工程流体力学考研分析》里有一句最关键的话（原文照录）：

\begin{quote}
\textbf{每年的真题大部分都源于课后习题的改编，课后习题一定要特别认真对待。}
\end{quote}

2026 年回忆版也印证了这一点：当年 5 道计算题里，
\textbf{第 3 题直接是教材动量方程习题 4.12、第 5 题直接是教材管嘴出流的例题}。
既然例题能原样进场，课后习题的改编题自然更多。
所以本册的定位是：\textbf{不是"补充练习"，而是真题的题库本体}。

\section*{二、本册怎么编排}

\begin{center}
\begin{tabularx}{\textwidth}{@{}clX@{}}
\toprule
部分 & 内容 & 用法 \\
\midrule
第一部分 & 第 1--8 章\textbf{全部课后习题}（题目 + 星级） & 先合上解答自己写，写不出再看提示 \\
第二部分 & \textbf{课后习题解答}（统一置于书末） & 对答案、看规范步骤与评分点 \\
\bottomrule
\end{tabularx}
\end{center}

\noindent\textbf{关于答案的来源：}赵琴教材课后习题\textbf{原书不附答案}。
本册解答由编者按教材正文的公式与例题方法\textbf{逐步推导}得到，
每一步都写清"依据哪条公式、为什么可以这样列式"，
并对易错处用红框标出。若某题条件不足或有多种解法，会在题后注明。

\section*{三、星级怎么来的（有文件依据，不是印象）}

\begin{center}
\begin{tabularx}{\textwidth}{@{}clX@{}}
\toprule
星级 & 含义 & 判定条件（满足其一） \\
\midrule
\xstarfull{5} & 必做 & 该考点在历年真题中作为计算题出现 $\ge 3$ 次；或考情分析明写"考试计算题重点" \\
\xstarfull{4} & 高频 & 真题中出现 2--3 次；或考情分析写"重点掌握" \\
\xstarfull{3} & 常考 & 真题中出现 1 次；或考点分析写"掌握基本概念与公式套用" \\
\xstarfull{2} & 了解 & 真题未出现，但属考纲核心内容 \\
\xstarfull{1} & 边缘 & 概念性/扩展题 \\
\bottomrule
\end{tabularx}
\end{center}

\noindent 每题的星级后都附〔依据：$\cdots$〕，写明是\textbf{哪一年的第几题}或\textbf{哪份分析文件的哪一页}。
完整的证据清单（E1--E6）与逐章星级总表，见《818 教材精讲》附录 A。

\section*{四、逐章题型分布（据教材目录整理）}

\begin{center}
\begin{tabularx}{\textwidth}{@{}clXc@{}}
\toprule
章 & 名称 & 习题范围（教材书页） & 题量 \\
\midrule
1 & 绪论 & 1.1--1.11（书 p8--9） & 11 \\
2 & 流体静力学 & 2.1--2.17（书 p24--27） & 17 \\
3 & 流体运动学 & 3.1--3.11（书 p44 起） & 11 \\
4 & 流体动力学基本方程 & 4.1--4.16（书 p64 起） & 16 \\
5 & 管路、孔口、管嘴的水力计算 & 5.1--5.20（书 p95--98） & 20 \\
6 & 相似理论与量纲分析 & 6.1--6.12（书 p116--118） & 12 \\
7 & 理想流体动力学 & 7.1--7.12（书 p150--151） & 12 \\
8 & 黏性流体动力学基础 & 8.1--8.16（书 p183--184） & 16 \\
\bottomrule
\end{tabularx}
\end{center}

\begin{yicuo}
\textbf{不要只做"看起来像真题"的题。}818 的选择题、判断题大量改编自课后习题里的\textbf{概念题}
（如 1.1 连续介质模型的必要性、1.3 黏性随温度的变化趋势、2.10 压力中心与形心重合的条件）。
这类题做起来快、收益高，务必逐题写下答案要点。
\end{yicuo}
